\chapter{{\fontsize{14pt}{21pt}\selectfont\MakeUppercase{Приложение Б. Программа и методика испытаний}}}
\label{app:pimi}

\section*{Б.1 Объект испытаний}

Объектом испытаний является программный комплекс корпоративного обучения: серверная часть на базе FastAPI, веб-панель на React и TypeScript, мобильное приложение на Flutter, СУБД PostgreSQL (или SQLite в режиме разработки), конфигурация Docker Compose.

\section*{Б.2 Цель испытаний}

Целью испытаний является подтверждение соответствия реализации функциональным и нефункциональным требованиям, описанным в техническом задании (приложение~А), на критическом пути пользователей и при типовых операциях администрирования.

\section*{Б.3 Условия проведения испытаний}

Испытания проводятся на стенде, развёрнутом по инструкции главы~3 пояснительной записки и приложению~Г, с использованием демонстрационных учётных записей и заполненной демонстрационной базы (\texttt{seed\_demo}).

\section*{Б.4 Средства испытаний}

\begin{enumerate}
  \item персональная ЭВМ администратора с браузером и доступом к веб-панели;
  \item мобильное устройство или эмулятор с установленным клиентским приложением;
  \item средства автоматизированного тестирования: Pytest (серверная часть), Vitest и Playwright (web), Flutter test (mobile) --- по стратегии репозитория \texttt{docs/testing-strategy.md}.
\end{enumerate}

\section*{Б.5 Методы испытаний}

Применяются функциональное тестирование по тест-кейсам, интеграционное тестирование API, обзорное тестирование пользовательского интерфейса, проверка ограничений доступа по ролям и идентификатору организации.\cite{iso29119,kaner_testing_en,gost1930179}

\section*{Б.6 Порядок проведения и отчётность}

Фиксируются версия сборки, дата прогона, результаты каждого тест-кейса (пройден / не пройден), идентификатор дефекта при отказе. Итоговый отчёт включает сводную таблицу и список замечаний.

\section*{Б.7 Таблица тест-кейсов критического пути}

\begin{table}[H]
\centering
\caption{Результаты испытаний по тест-кейсам критического пути}
\label{tab:appb-tests}
{\setstretch{1}\fontsize{11pt}{13pt}\selectfont
\setlength{\tabcolsep}{3pt}
\begin{tabularx}{\textwidth}{|>{\centering\arraybackslash}p{10mm}|>{\RaggedRight\arraybackslash}p{38mm}|>{\RaggedRight\arraybackslash}X|>{\RaggedRight\arraybackslash}p{28mm}|>{\centering\arraybackslash}p{18mm}|}
\hline
\textbf{ID} & \textbf{Описание} & \textbf{Шаги и ожидаемый результат} & \textbf{Фактический результат} & \textbf{Статус} \\
\hline
TC-01 & Регистрация пользователя & POST \texttt{/auth/register}, создание записи пользователя & создан & ОК \\
\hline
TC-02 & Вход и получение JWT & POST \texttt{/auth/login}, выдача access и refresh токенов & выданы & ОК \\
\hline
TC-03 & Выбор организации & POST \texttt{/tenants/select}, активное членство & организация выбрана & ОК \\
\hline
TC-04 & Создание курса & POST \texttt{/courses}, запись с \texttt{tenant\_id} & курс создан & ОК \\
\hline
TC-05 & Создание урока & POST \texttt{/lessons}, связь с курсом & урок создан & ОК \\
\hline
TC-06 & Создание теста и вопросов & POST \texttt{/tests}, POST вопросов & тест доступен & ОК \\
\hline
TC-07 & Назначение курса & POST назначения группе или пользователю & назначение создано & ОК \\
\hline
TC-08 & Список курсов слушателя & GET курсов под ролью learner & список возвращён & ОК \\
\hline
TC-09 & Старт попытки теста & POST \texttt{/tests/\{id\}/start} & попытка создана & ОК \\
\hline
TC-10 & Следующий вопрос & GET \texttt{/attempts/\{id\}/next-question} & вопрос или завершение & ОК \\
\hline
TC-11 & Отправка ответа & POST \texttt{/attempts/\{id\}/submit-answer} & обратная связь по сложности & ОК \\
\hline
TC-12 & Завершение попытки & POST \texttt{/attempts/\{id\}/finish} & результат и рекомендации & ОК \\
\hline
TC-13 & Рекомендации слушателя & GET \texttt{/recommendations/me} & список тем & ОК \\
\hline
TC-14 & Аналитика дашборда & GET \texttt{/analytics/dashboard} & агрегаты & ОК \\
\hline
TC-15 & Изоляция арендаторов & запрос данных другого \texttt{tenant\_id} под чужой организацией & отказ 403/404 & ОК \\
\hline
TC-16 & Доступ без токена & запрос к защищённому эндпоинту & 401 & ОК \\
\hline
TC-17 & Загрузка медиа & POST \texttt{/media} с файлом & URL медиа & ОК \\
\hline
TC-18 & Прогресс урока & POST прогресса страницы урока & сохранено & ОК \\
\hline
TC-19 & Веб-панель вход & страница авторизации и переход в shell & отображение разделов & ОК \\
\hline
TC-20 & Мобильный поток теста & экран TestFlowScreen: старт, ответы, финиш & завершено без падений & ОК \\
\hline
\end{tabularx}
}
\end{table}

\noindent\textit{Примечание: колонки «Фактический результат» и «Статус» содержат ориентировочные значения для демонстрационной сборки; при защите подставляются фактические результаты прогона.}
